\section{C-Learner Has Finite Variance in a Simple Low-Overlap Setting, While AIPW and TMLE Do Not}
\label{sec:theory_simple}
We have discussed how the direct method (naive plug-in estimator) is stable, including in challenging settings with low overlap, but lacks the asymptotic properties of debiased methods like one-step estimation and targeting. On the other hand, basic versions of one-step estimation and targeting achieve asymptotic optimality but are unstable in settings with low overlap. It appears C-Learner is the best of both worlds, as it has desirable asymptotic properties (\Cref{sec:theory_text}) while producing stable estimates (\Cref{sec:experiments}). 
Here, we use a simple theoretical example to show how C-Learner inherits the stability of the direct method. 
As C-Learner, AIPW, and TMLE are all asymptotically equivalent under (roughly) standard assumptions (\Cref{sec:theory_text}), the assumptions in this section are necessarily nonstandard. 
We do not claim that C-Learner is strictly superior to AIPW or TMLE; rather, we articulate specific scenarios in which C-Learner outperforms AIPW and TMLE. 


The intuition is that we have outcome models that are well-behaved, but inverse propensity terms (such as $A/\what\pi(X)\cdot (Y-\what\mu(X))$) that are very poorly behaved, so that estimators that are plug-ins of outcome models have finite variance but estimators with additive inverse propensity terms have infinite variance. Such poorly-behaved inverse propensity terms can happen when there is low overlap. 


More specifically, we can specify the function class of outcome models $\mathcal F$ to only consist of outcome models contained within some square-integrable envelope $C(x)$ with $P[C(X)^2]<\infty$, so that $P_{k,n}[f(X)]$ for any $f\in \mathcal F$ has finite variance. 
At the same time, we make assumptions so that the inverse propensity term in the AIPW and TMLE have infinite variance, for example, because $1/\pi(X)$ is heavy-tailed and not integrable due to low overlap. 
For simplicity, we also assume the true propensity weights are known and used for estimation. 
Under these assumptions, the AIPW and TMLE estimators do not have finite variance, while C-Learner does have finite variance. 
Note that the assumption of heavy-tailed $1/\pi(X)$ is a departure from standard assumptions in Section~\ref{sec:theory_text}, where we assume that $\pi(X)$ is bounded away from 0. 




Although assuming known propensity scores is quite strong, the results are still interesting as AIPW, TMLE, and C-Learner all use these propensity weights but only C-Learner has finite variance. 
Additionally, although our example will show that the direct method has finite variance like the C-Learner, our goal is to compare with estimators that are asymptotically optimal under more standard assumptions (\Cref{sec:theory_text}).

\subsection{Setting}
We focus on the mean missing outcome setting, as described in \Cref{sec:ate}.
Consider binary treatments (actions) $A\in\{0,1\}$, covariates $X\in\R^d$, and potential outcomes $Y(1)$ and $Y(0):=0$. Let $(X,A,Y)$ be i.i.d. observations drawn from distribution $P$. Our goal is to estimate the mean missing outcome $\psi:=\int Y(1) dP$, which we assume is finite.  




Let $\|\cdot \|$ denote the Euclidean norm. For consistency with previous sections, let $P[\cdot]$ denote the expectation operator under $P$. 
We consider the cross-fitted~\citep{van2011cross,ChernozhukovChDeDuHaNeRo18} estimators as defined in \Cref{sec:data_splitting}, but using true propensities $\pi$ instead of fitted ones $\what\pi_{-k,n}$. 



\subsection{Assumptions}
To show that C-Learner remains well-behaved even when inverse propensity terms blow up, we assume all outcome models lie within a square-integrable envelope $C(x)$, %
i.e. for any $f\in \mathcal F$,
$|f(X)|\leq C(X)$ almost surely, and $P[C(X)^2]<\infty$. %
This is much weaker than assuming that the model class is bounded everywhere, for example. Below we provide examples, such as where $C(x)$ is a maximum of linear functions of $x$. Note that bounding a function class with square-integrable envelopes is also standard in empirical process theory. 

We note that while our proposed methodology for C-Learner (\Cref{sec:methodology})  does not constrain outcome models to be within any specific bounded envelopes, one can choose a wide enough envelope to encompass reasonable outcome models for the setting, and then proceed with C-Learner within this newly constrained function class (i.e. let $\mathcal F$ be outcome models constrained to be within the chosen envelope). 

\begin{assumption}[Outcome models are contained within a square-integrable envelope]
\label{ass:envelope}
There is a square-integrable envelope $C(x)$ (i.e. $P[C(X)^2]<\infty$) such that all elements in the model class $f\in \mathcal F$ are contained in the envelope (i.e. $|f(X)|\leq C(X)$ almost surely).\end{assumption}
Below we give examples to show that a linear function is square-integrable if $X$ has finite second moments, and that the maximum of square-integrable functions is also square-integrable. Thus, if we assume $X$ has finite second moments, then one can construct square-integrable envelopes $C(x)$ that are the maximum of linear functions;
this shows that assuming $f$ is contained in a square-integrable envelope is a weaker assumption than assuming $f$ is bounded everywhere by constants, for example.

\begin{example}[Linear models with finite coefficients are square-integrable if $X$ has finite second moments]
\label{ex:linear_square_int}
    Let $g(x):=\beta^\top x$ for some fixed $\beta$. Assume $P[\|X\|^2]<\infty$. Then $g(x)$ is square-integrable. %
    This follows by Cauchy-Schwarz: 
    $P[(\beta^\top X)^2]\leq \|\beta\|^2 P[\|X\|^2]<\infty.$
\end{example}


\begin{example}[Maximum of a finite number of square-integrable functions is square-integrable]\label{ex:max_square_int}
Consider $g_1,\ldots,g_k$ that are all square-integrable so that $P[g_i(X)^2]<\infty$ for $i=1,\ldots,k$. 
Note that $\max(g_1(x)^2,\ldots,g_k(x)^2)\leq |g_1(x)|^2+\ldots+ |g_k(x)|^2$
so that taking expectations %
$$P[\max(|g_1(X)|,\ldots,|g_k(X)|)^2]
= P[\max(g_1(X)^2,\ldots,g_k(X)^2)]
\leq \sum_{i=1}^k P[|g_i(X)|^2]<\infty$$
with the last inequality by square-integrability of each $g_i$. 
\end{example}

Below, we also assume that a feasible solution to the constraint exists in $\mathcal F$ with high probability. This high probability depends on the size of the envelope. In Example~\ref{ex:feasible_envelope_whp} we present a simple example where for a fixed $\epsilon>0$, we can choose $\mathcal F$ (and envelope $C(x)$) such that the feasible solution to the constraint exists in $\mathcal F$  with probability $\geq 1-\epsilon$. 

\begin{assumption}[C-Learner solution exists in $\mathcal F$]
\label{ass:soln_exists}
With probability $\geq 1-\epsilon$ over draws of data of size $n$, 
the C-Learner solution exists in $\mathcal F$. \end{assumption}
A simple sufficient (but not necessary) condition for the C-Learner solution to exist in $\mathcal F$ is that the constant functions corresponding to the maximum and minimum values in $P_{k,n}$, which we denote $f_+$ and $f_-$, respectively, are both in $\mathcal F$, and $\mathcal F$ is closed under convex combinations. This is sufficient because 
$P_{k,n}\left[\frac{A}{\pi(X)}(Y-f_+(X))\right]\leq 0$ and $P_{k,n}\left[\frac{A}{\pi(X)}(Y-f_-(X))\right]\geq 0$, 
so there will be a convex combination of $f_+$ and $f_-$ for which the constraint is exactly 0. In Example~\ref{ex:feasible_envelope_whp} below we construct a model class $\mathcal F$ contained in envelope $C(x)$ where $\mathcal F$ and $C(x)$ are defined in such a way that the maximum and minimum values of $Y$ in the sample $P_{k,n}$ are within $\mathcal F$ and envelope $C(x)$ at least $1-\epsilon$ of the time. Note that randomness over $P_{-k,n}$ is not directly relevant to this condition. 

\begin{example}[Constructing $\mathcal F$ where C-Learner solution exists in $\mathcal F$ with high probability]
\label{ex:feasible_envelope_whp}
Consider a model class $\mathcal F_0$ that contains functions of choice, and also all constant functions, that is closed under convex combinations. Also consider an envelope function of choice $C_0(x)$ that is chosen to contain most reasonable unconstrained solutions. There are many choices for $C_0$ (\Cref{ex:linear_square_int},~\Cref{ex:max_square_int}). It is not required that $|f(X)|\leq C_0(X)$ almost surely for all $f\in \mathcal F_0$. 
Fix $m$, the number of samples in $P_{k,n}$, and probability threshold $\epsilon>0$. 
Assume $Y$ is continuous and that the CDF of $Y$ is known: $F(x)=P(Y\leq x)$. 
Let $\alpha=\frac{1}{2}(1-(1-\epsilon)^{1/m})$. 
Then consider lower and upper bounds
$a=F^{-1}(\alpha), \; b=F^{-1}(1-\alpha).$
By construction of $a$ and $b$, 
$P( \cap_{i=1}^m \{Y_i\in[a,b]\})\geq 1-\epsilon$
so that we can choose a new envelope
$$C(x)=\max(C_0(x), |a|, |b|)$$
that contains each of $Y_1,\ldots,Y_m$ with probability $\geq 1-\epsilon$ and is square-integrable (\Cref{ex:max_square_int}). 
Then define $\mathcal F$ to be all $f\in\mathcal F_0$ for which $|f(x)|\leq C(x)$ almost surely.
\end{example}

Lastly, we make the following assumption so that the inverse propensity terms in AIPW and TMLE have infinite variance. We state the assumption, then discuss its two parts. 
\begin{assumption}[Inverse propensity term is not square-integrable]
\label{ass:prop_not_square_int}
(I) There is some $v_0>0$ such that 
for all $X$ in the support of $P$, $\text{Var}(Y\mid X, A=1)\geq v_0.$ (II) Also, $\pi(X)>0$ almost surely and
$P\left[\frac{1}{ \pi(X)}\right]=\infty.$
\end{assumption}
In the first part of the assumption above, we assume that the conditional variance of the outcome, conditional on $X$ and treatment $A=1$, is above a threshold $v_0>0$. In the vast majority of real settings, the conditional variance of the outcome will be nonzero, and it may be plausible to assume a lower bound.

In the second part of the assumption above, we require that 
$P[1/\pi(X)]=\infty$, which corresponds to low overlap settings.
We also assume $\pi(X)>0$ almost surely so the estimators are well-defined. %
There are many choices of $\pi(X)$ for which $P[1/\pi(X)]=\infty$; we provide simple examples. %


\begin{example}[Beta $\pi(X)$]
Let $\pi(X)\sim \text{Beta}(a,1)$ with density
$f(p)=a\,p^{a-1}$, then $P[1/\pi(X)]=a/(a-1)$ if $a>1$, and $P[1/\pi(X)]=\infty$ otherwise. 
As a simple example, if $\pi(X)\sim\text{Unif}(0,1)$ then $\pi(X)\sim\text{Beta}(1,1)$, and $P[1/\pi(X)]=\infty$. 
\end{example}
\subsection{Result}
Now we proceed to the main result, which is that in the specific theoretical scenario defined above, C-Learner has finite variance, while AIPW and TMLE do not. 
\begin{theorem}
\label{thm:simple_variance}
Assume Assumptions~\ref{ass:envelope},~\ref{ass:soln_exists},~\ref{ass:prop_not_square_int}, 
i.e. fix an $\epsilon>0$, assume true propensities $\pi(x)$ are known, and assume the following:
\begin{enumerate}
    \item There is an envelope $C(x)$ with  $P[C(X)^2]<\infty$ such that for all $f\in \mathcal F$, $|f(X)|\leq C(X)$ almost surely. 
    \item With probability $\geq 1-\epsilon$ over draws of data of size $n$,  C-Learner solution exists in $\mathcal F$.
\item There is a $v_0>0$ such that 
$\text{Var}(Y\mid X, A=1)\geq v_0$
for all $X$ in the support of $P$. Also, $\pi(X)>0$ almost surely and
$P\left[\frac{1}{ \pi(X)}\right]=\infty$.
\end{enumerate}
Then the TMLE and AIPW estimators have infinite variance. In contrast, 
with probability $\geq 1-\epsilon$,
the C-Learner estimator has finite variance. 
\end{theorem}
See \Cref{sec:theory_simple_more} for the proof of \Cref{thm:simple_variance}. 

\label{sec:theory_simple-end}
